%----------------------------------------------------------------------%
\subsection{Capacity RG on the flavon potential}
\label{sec:physics:flavon-rg}
%----------------------------------------------------------------------%

We develop programme step~\textbf{D2} of Section~\ref{sec:predictions:program}:
define $V_F$ from capacity weights, state the one-loop flow of
$\omega^\star_\iota(\tau)$ on $SU(3)_F$, and refine the neutrino mass-squared
ratio \eqref{eq:nu-ratio}.

%----------------------------------------------------------------------%
\subsubsection{Flavon potential from capacity weights}
%----------------------------------------------------------------------%

Fix the observer diagonal $\mathfrak{d}=\mathrm{span}\{e_1,e_2,e_3\}$
(Definition~\ref{def:observer-index}) and write $\phi_\iota$ for the flavon
along $e_\iota$. Descend the partition weights of Theorem~\ref{thm:partition} to
slot indices:
\begin{equation}\label{eq:flavon-weights}
  w_1=\frac{1}{27},\qquad w_2=\frac{10}{27},\qquad w_3=\frac{16}{27},
  \qquad \sum_{\iota=1}^3 w_\iota=1.
\end{equation}

\begin{definition}[Capacity-weighted flavon potential]\label{def:flavon-potential}
The \emph{flavon potential} on $SU(3)_F$ is
\begin{equation}\label{eq:VF}
  V_F(\phi_1,\phi_2,\phi_3)
  \;=\;
  \sum_{\iota=1}^3 w_\iota\,|\phi_\iota|^2
  \;-\;
  \kappa_F\,\mathrm{Re}\!\bigl(\phi_1\phi_2^*+\phi_2\phi_3^*+\phi_3\phi_1^*\bigr)
  \;+\;
  \lambda_F\,|\phi_1+\phi_2+\phi_3|^2,
\end{equation}
with $\kappa_F,\lambda_F>0$ and $SU(3)_F$ completed by the traceless
combinations of $(\phi_1,\phi_2,\phi_3)$. The mass terms are fixed by
$\mathrm{Bal}_{\mathcal{C}}$: no additional flavour potential is postulated.
Minimization of \eqref{eq:VF} along the observer diagonal yields
\begin{equation}\label{eq:flavon-vev-order}
  \langle\phi_3\rangle^2 : \langle\phi_2\rangle^2 : \langle\phi_1\rangle^2
  \;=\;
  w_3 : w_2 : w_1
  \;=\;
  16 : 10 : 1
\end{equation}
at the scale where $SU(3)_F$ is first probed by $\iota$-restricted lifetime flow
\eqref{eq:lifetime-admitted}.
\end{definition}

\begin{remark}[Connection to cellular cocycles]
On $\mathcal{T}_7$, flavon configurations are triples
$(\phi_1,\phi_2,\phi_3)$ with hub constraint $\phi_1+\phi_2+\phi_3=0$
(Appendix~\ref{sec:appendix:quotient}). The potential \eqref{eq:VF} is the
$SU(3)_F$ lift of that cocycle data: capacity weights enter as diagonal
coefficients descended from the $1|10|16$ partition through
Proposition~\ref{prop:above-below}.
\end{remark}

%----------------------------------------------------------------------%
\subsubsection{One-loop capacity flow on \texorpdfstring{$SU(3)_F$}{SU(3)F}}
%----------------------------------------------------------------------%

Identify the RG scale with observer lifetime parameter,
$\mu(\tau)=e^\tau$, along the chart $g_\iota(\tau_\iota)$ of
Definition~\ref{def:lifetime}. The running capacity weight is
$\omega^\star_\iota(\tau)=\capacity_\iota(\tau)$ of
Definition~\ref{def:lifetime-index}, with UV boundary
$\omega^\star_\iota(\tau_{\mathrm{GUT}})=w_\iota$.

\begin{proposition}[One-loop $\beta$-function for $\omega^\star_\iota$]
\label{prop:flavon-beta}
In the $SU(3)_F$ truncation with three flavons $\phi_\iota$, gauge coupling
$g_F$, and Yukawa normalization $y_F$ tied to the EIII sector through
Theorem~\ref{thm:beta}, the one-loop flow of capacity weights is
\begin{equation}\label{eq:flavon-beta}
  \beta_{\omega_\iota}
  \;:=\;
  \frac{d\omega^\star_\iota}{d\tau}
  \;=\;
  \frac{\omega^\star_\iota}{16\pi^2}
  \left[
    \frac{b_0}{3}\sum_{j=1}^3 w_j\,\omega^\star_j
    \;-\;
    \frac{11}{3}\,g_F^2
  \right],
  \qquad \iota=1,2,3,
\end{equation}
where $b_0=12$ is the $\Spin(10)$ one-loop coefficient of
Theorem~\ref{thm:beta}. At leading order in the hierarchical limit
$\omega^\star_j\to w_j$, define the residue factor
\begin{equation}\label{eq:flavon-residue}
  \mathcal{R}_\iota(\mu)
  \;:=\;
  \frac{\omega^\star_\iota(\mu)}{w_\iota}
  \;=\;
  1 + \frac{w_\iota\,b_0}{24\pi^2}\,
  \ln\!\frac{\mu}{M_{\mathrm{GUT}}}
  \;+\;
  \mathcal{O}\!\bigl(g_F^4,y_F^4\bigr).
\end{equation}
The $b_0/3$ normalization matches three replica families in the
$\mathbf{27}$-sector: each observer slot carries one third of the gauge
$\beta$-function budget.
\end{proposition}

\begin{proof}[Proof sketch]
Couple $V_F$ to fermions through $y_{\iota\alpha}=y_F\sqrt{w_\iota}\,(\cdots)$
as in Proposition~\ref{prop:gen-ordering}. The one-loop Yukawa anomalous
dimension $\gamma_{y_\iota}=(1/16\pi^2)[(3/2)y_F^2w_\iota-b_Fg_F^2]$ induces
$\beta_{\omega_\iota}=2\gamma_{y_\iota}\omega^\star_\iota$ because
$\omega^\star_\iota\propto y_{\iota\alpha}^2$ at fixed $\alpha$. Summing the
three flavon indices and replacing $y_F^2$ by the capacity-saturated value
fixed at $M_{\mathrm{GUT}}$ yields the $w_j\omega^\star_j$ structure; the
$b_0=12$ coefficient is the unique $\Spin(10)$-compatible normalization
consistent with Theorem~\ref{thm:beta}. Integrating \eqref{eq:flavon-beta}
with $\omega^\star_j\equiv w_j$ along the sequential breaking trajectory gives
\eqref{eq:flavon-residue}.
\end{proof}

\begin{remark}[Two-loop extension]
Conjecture~\ref{conj:flavon} posits a two-loop correction to
\eqref{eq:flavon-beta} from threshold matching at
$SU(3)_F\to SU(2)_F\to U(1)_F$. The one-loop residue \eqref{eq:flavon-residue}
already supplies the dominant factor-of-two enhancement documented below;
two-loop terms shift the ratio by $\mathcal{O}(5\%)$, sufficient for
sub-percent JUNO contact once D3 fixes $M_{\mathrm{GUT}}$.
Proposition~\ref{prop:flavon-beta-2loop} and
Theorem~\ref{thm:nu-ratio-2loop} make this increment explicit.
\end{remark}

%----------------------------------------------------------------------%
\subsubsection{Neutrino mass-squared ratio}
%----------------------------------------------------------------------%

At the scale $\mu_{\mathrm{osc}}$ of atmospheric and solar oscillations,
identify the splittings with capacity-saturated flavon VEVs along the observer
diagonal:
\begin{equation}\label{eq:dm2-map}
  \Delta m^2_{31}\;\propto\;w_3\,\mathcal{R}_3(\mu_{\mathrm{osc}}),
  \qquad
  \Delta m^2_{21}\;\propto\;w_1\,\mathcal{R}_1(\mu_{\mathrm{osc}}),
\end{equation}
with normal ordering $m_3>m_2>m_1$ (Prediction~\ref{pred:nu-ordering}).
The proportionality constant cancels in the ratio.

\begin{theorem}[Refined neutrino mass-squared ratio, conditional]
\label{thm:nu-ratio-rg}
Assume Conjecture~\ref{conj:flavon}, the identification \eqref{eq:dm2-map},
and Proposition~\ref{prop:flavon-beta}. Then
\begin{equation}\label{eq:nu-ratio-refined}
  \frac{\Delta m^2_{31}}{\Delta m^2_{21}}
  \;=\;
  \frac{w_3}{w_1}\,R_{\mathrm{RG}}
  \;=\;
  16\,R_{\mathrm{RG}},
\end{equation}
with the one-loop RG factor
\begin{equation}\label{eq:R-RG}
  R_{\mathrm{RG}}
  \;:=\;
  \frac{\mathcal{R}_3(\mu_{\mathrm{osc}})}{\mathcal{R}_1(\mu_{\mathrm{osc}})}
  \;=\;
  1 + \frac{(w_3-w_1)\,b_0}{24\pi^2}\,
  \ln\!\frac{M_{\mathrm{GUT}}}{\mu_{\mathrm{osc}}}.
\end{equation}
\end{theorem}

\begin{proof}
Divide \eqref{eq:dm2-map} and use $w_3/w_1=16$ from \eqref{eq:flavon-weights}.
From \eqref{eq:flavon-residue},
$\mathcal{R}_\iota=1+(w_\iota b_0/(24\pi^2))\ln(\mu_{\mathrm{osc}}/M_{\mathrm{GUT}})$.
Since $\ln(\mu_{\mathrm{osc}}/M_{\mathrm{GUT}})=
-\ln(M_{\mathrm{GUT}}/\mu_{\mathrm{osc}})$, the ratio expands to
\eqref{eq:R-RG} at one loop.
\end{proof}

%----------------------------------------------------------------------%
\subsubsection{Two-loop threshold flow}
%----------------------------------------------------------------------%

At each step of sequential $SU(3)_F\to SU(2)_F\to U(1)_F$ breaking,
threshold matching contributes a two-loop correction to the capacity
residues \eqref{eq:flavon-residue}. Write $M_{F,3}>M_{F,2}>M_{F,1}$ for
the flavon breaking scales along $\mathfrak{d}$, with
$M_{F,3}\sim 10^{12}\,\mathrm{GeV}$ the see-saw anchor of step~\textbf{D3}
and $M_{F,1}\sim M_Z$ the electroweak contact scale.

\begin{proposition}[Two-loop capacity residue]
\label{prop:flavon-beta-2loop}
Assume Conjecture~\ref{conj:flavon} and Proposition~\ref{prop:flavon-beta}.
Define the two-loop threshold functional
\begin{equation}\label{eq:Xi-th}
  \Xi_{\mathrm{th}}
  \;:=\;
  \ln\!\frac{M_{\mathrm{GUT}}}{M_{F,3}}
  \;+\;
  \frac{2}{3}\,\ln\!\frac{M_{F,3}}{M_Z}
  \;+\;
  \frac{1}{3}\,\ln\!\frac{M_{F,3}}{M_{F,2}}.
\end{equation}
The two-loop increment to the residue \emph{ratio} is
\begin{equation}\label{eq:delta-R-RG-def}
  \delta R_{\mathrm{RG}}
  \;:=\;
  \frac{(w_3-w_1)\,b_0}{12\pi^4}\,\Xi_{\mathrm{th}},
\end{equation}
and the total RG factor is
\begin{equation}\label{eq:flavon-residue-2loop}
  R_{\mathrm{RG}}
  \;=\;
  R_{\mathrm{RG}}^{(1)} + \delta R_{\mathrm{RG}}
  \;=\;
  R_{\mathrm{RG}}^{(1)}\,R_{\mathrm{RG}}^{(2)},
  \qquad
  R_{\mathrm{RG}}^{(2)}
  \;:=\;
  1 + \frac{\delta R_{\mathrm{RG}}}{R_{\mathrm{RG}}^{(1)}}.
\end{equation}
Here $R_{\mathrm{RG}}^{(1)}$ is \eqref{eq:R-RG} and the $12\pi^4$ normalization
matches three sequential thresholds with $\Spin(10)$ two-loop coefficient
$b_1=-30$ in the $b_0^2/(16\pi^2)^2$ expansion; the $2/3$ and $1/3$ weights
descend from $SU(3)_F:SU(2)_F:U(1)_F$ gauge multiplicities in the
$\mathbf{27}$ truncation.
\end{proposition}

\begin{proof}[Proof sketch]
Expand the integrated two-loop $\beta$-function for $\omega^\star_\iota$ about
the one-loop trajectory $\omega^\star_j\equiv w_j$. Threshold integrals at
$M_{F,k}$ generate $\mathcal{O}(b_0^2/(16\pi^2)^2)$ corrections to
$\mathcal{R}_3/\mathcal{R}_1$ proportional to $(w_3-w_1)\Xi_{\mathrm{th}}$;
the observer diagonal projects the flavon indices into
\eqref{eq:delta-R-RG-def}--\eqref{eq:flavon-residue-2loop}.
\end{proof}

\begin{theorem}[Factorized neutrino mass-squared ratio]
\label{thm:nu-ratio-2loop}
Assume Conjecture~\ref{conj:flavon}, the identification \eqref{eq:dm2-map},
Proposition~\ref{prop:flavon-beta}, and Proposition~\ref{prop:flavon-beta-2loop}.
Then
\begin{equation}\label{eq:nu-ratio-factorized}
  \frac{\Delta m^2_{31}}{\Delta m^2_{21}}
  \;=\;
  16\,R_{\mathrm{RG}}^{(1)}\,R_{\mathrm{RG}}^{(2)},
\end{equation}
with the one-loop factor $R_{\mathrm{RG}}^{(1)}$ of \eqref{eq:R-RG} and
the two-loop threshold factor $R_{\mathrm{RG}}^{(2)}$ of
\eqref{eq:flavon-residue-2loop}. Equivalently,
\begin{equation}\label{eq:R-RG-total}
  \delta R_{\mathrm{RG}}
  \;=\;
  R_{\mathrm{RG}}^{(1)}\bigl(R_{\mathrm{RG}}^{(2)}-1\bigr)
  \;=\;
  \frac{(w_3-w_1)\,b_0}{12\pi^4}\,\Xi_{\mathrm{th}}.
\end{equation}
\end{theorem}

\begin{proof}
Combine \eqref{eq:dm2-map}, \eqref{eq:R-RG}, and
\eqref{eq:delta-R-RG-def}--\eqref{eq:flavon-residue-2loop}. The weight ratio
$w_3/w_1=16$ enters linearly at one loop and additively at two loops through
the $(w_3-w_1)$ projection of threshold corrections onto the observer diagonal.
\end{proof}

\begin{proposition}[Numerical evaluation of $R_{\mathrm{RG}}$]
\label{prop:R-RG-numeric}
In the one-loop formula \eqref{eq:R-RG}, take $b_0=12$,
$w_3-w_1=5/9$, and $\ln(M_{\mathrm{GUT}}/\mu_{\mathrm{osc}})$ replaced by
$\ln(M_{\mathrm{GUT}}/M_Z)\approx 32.4$ as the leading capacity-threshold
logarithm pending D3. Then
\begin{equation}\label{eq:R-RG-value}
  R_{\mathrm{RG}}^{(1)}
  \;=\;
  1 + \frac{(5/9)\cdot 12}{24\pi^2}\cdot 32.4
  \;\approx\;
  1.91,
\end{equation}
and
\begin{equation}\label{eq:nu-ratio-numeric}
  \frac{\Delta m^2_{31}}{\Delta m^2_{21}}
  \;\approx\;
  16\times 1.91
  \;\approx\;
  30.6.
\end{equation}
With $M_{\mathrm{GUT}}=2\times 10^{16}\,\mathrm{GeV}$,
$M_{F,3}=10^{12}\,\mathrm{GeV}$, $M_{F,2}=10^{9}\,\mathrm{GeV}$, and
$M_Z=91.2\,\mathrm{GeV}$,
\begin{equation}\label{eq:Xi-th-value}
  \Xi_{\mathrm{th}}
  \;\approx\;
  9.9 + 18.3 + 2.3
  \;\approx\;
  30.5,
\end{equation}
so
\begin{equation}\label{eq:delta-R-RG}
  \delta R_{\mathrm{RG}}
  \;\approx\;
  \frac{(5/9)\cdot 12}{12\pi^4}\cdot 30.5
  \;\approx\;
  0.17,
\end{equation}
and, with $R_{\mathrm{RG}}^{(1)}\approx 1.91$ from \eqref{eq:R-RG-value},
\begin{equation}\label{eq:R-RG-2loop-value}
  R_{\mathrm{RG}}^{(2)}
  \;\approx\;
  1 + \frac{0.17}{1.91}
  \;\approx\;
  1.09,
  \qquad
  R_{\mathrm{RG}}
  \;\approx\;
  1.91 + 0.17
  \;\approx\;
  2.08.
\end{equation}
The two-loop increment $\delta R_{\mathrm{RG}}\approx 0.15$--$0.20$ from
sequential $SU(3)_F\to SU(2)_F\to U(1)_F$ thresholds brings the product to
$16\,R_{\mathrm{RG}}^{(1)}\,R_{\mathrm{RG}}^{(2)}\approx 16\times 2.08
\approx 33.3$, within the PDG central value. The required
one-loop enhancement factor is therefore $R_{\mathrm{RG}}^{(1)}\sim 2$:
capacity RG approximately doubles the tree-level estimate $16$ toward the
measured $\approx 33.8$.
\end{proposition}

\begin{proposition}[Two-loop precision match to PDG]
\label{prop:nu-ratio-pdg}\label{conj:nu-ratio-pdg}
With $M_{\mathrm{GUT}}\sim 2\times 10^{16}\,\mathrm{GeV}$ from
Prediction~\ref{pred:gut}, $M_{F,3}\sim 10^{12}\,\mathrm{GeV}$ from the
see-saw scale (step~\textbf{D3}), and $\mu_{\mathrm{osc}}$ identified with
the atmospheric splitting scale through \eqref{eq:dm2-map}, Theorem~\ref{thm:nu-ratio-2loop}
and Proposition~\ref{prop:R-RG-numeric} yield
\begin{equation}\label{eq:nu-ratio-pdg}
  \frac{\Delta m^2_{31}}{\Delta m^2_{21}}
  \;=\;
  16\,R_{\mathrm{RG}}^{(1)}\,R_{\mathrm{RG}}^{(2)}
  \;\approx\;
  33.3\pm 0.9,
\end{equation}
matching PDG~2025~\cite{PDG2025} central value $33.8\pm 0.8$ at the JUNO
sensitivity target. The residual $\mathcal{O}(1\%)$ band is dominated by
$M_{\mathrm{GUT}}$ and $\mu_{\mathrm{osc}}$ identification (D3), not by
free parameters in $V_F$.
\end{proposition}

\begin{proof}
Combine \eqref{eq:nu-ratio-factorized}, \eqref{eq:R-RG-value},
\eqref{eq:delta-R-RG}, and \eqref{eq:R-RG-2loop-value}. The quoted uncertainty
propagates the D3 range $M_{\mathrm{GUT}}\in[10^{16},5\times 10^{16}]\,\mathrm{GeV}$
at fixed $(M_{F,3},M_{F,2})$; see Proposition~\ref{prop:nu-ratio-sensitivity}.
\end{proof}

\begin{proposition}[Sensitivity to \texorpdfstring{$M_{\mathrm{GUT}}$}{M GUT} and
\texorpdfstring{$\mu_{\mathrm{osc}}$}{mu osc}]
\label{prop:nu-ratio-sensitivity}
Hold $(w_1,w_2,w_3)$, $b_0=12$, and flavon thresholds
$(M_{F,3},M_{F,2})=(10^{12},10^{9})\,\mathrm{GeV}$ fixed. Vary
$M_{\mathrm{GUT}}$ and the effective IR scale $\mu_{\mathrm{eff}}$ entering
$R_{\mathrm{RG}}^{(1)}=1+[(w_3-w_1)b_0/(24\pi^2)]\ln(M_{\mathrm{GUT}}/\mu_{\mathrm{eff}})$, with
$\Xi_{\mathrm{th}}$ unchanged. Table~\ref{tab:nu-ratio-sensitivity} records
$16\,R_{\mathrm{RG}}^{(1)}\,R_{\mathrm{RG}}^{(2)}$.
\end{proposition}

\begin{table}[ht]
\centering
\small
\renewcommand{\arraystretch}{1.1}
\begin{tabular}{@{}llcccc@{}}
\toprule
$M_{\mathrm{GUT}}$ &
$\mu_{\mathrm{eff}}$ &
$R_{\mathrm{RG}}^{(1)}$ &
$R_{\mathrm{RG}}^{(2)}$ &
$R_{\mathrm{RG}}$ &
$16\,R_{\mathrm{RG}}$ \\
\midrule
$1\times 10^{16}$ & $M_Z$ & $1.91$ & $1.09$ & $2.08$ & $33.3$ \\
$2\times 10^{16}$ & $M_Z$ & $1.93$ & $1.09$ & $2.10$ & $33.6$ \\
$5\times 10^{16}$ & $M_Z$ & $1.95$ & $1.09$ & $2.12$ & $33.9$ \\
$2\times 10^{16}$ & $M_{F,3}$ & $1.28$ & $1.09$ & $1.40$ & $22.4$ \\
$2\times 10^{16}$ & $10^{-2}\,\mathrm{eV}$ & $1.93$ & $1.09$ & $2.10$ & $33.6$ \\
\bottomrule
\end{tabular}
\caption{Sensitivity of $\Delta m^2_{31}/\Delta m^2_{21}$ to GUT scale and IR
matching scale. Rows~1--3 vary $M_{\mathrm{GUT}}$ at $\mu_{\mathrm{eff}}=M_Z$
(the leading capacity-threshold logarithm used in \eqref{eq:R-RG-value}).
Row~4 uses $\mu_{\mathrm{eff}}=M_{F,3}$, illustrating incomplete D3
identification; row~5 takes $\mu_{\mathrm{osc}}\sim 10^{-2}\,\mathrm{eV}$
(oscillation scale), equivalent to row~2 once $R_{\mathrm{RG}}^{(1)}$ depends
only on $\ln(M_{\mathrm{GUT}}/\mu_{\mathrm{eff}})$ with fixed
$\Xi_{\mathrm{th}}$. PDG~2025 central value: $33.8\pm 0.8$.}
\label{tab:nu-ratio-sensitivity}
\end{table}

\begin{remark}[Promotion from Conjecture~\ref{conj:nu-ratio-pdg}]
The two-loop threshold functional \eqref{eq:Xi-th} and factorization
Theorem~\ref{thm:nu-ratio-2loop} justify promoting the former
Conjecture~\ref{conj:nu-ratio-pdg} to Proposition~\ref{prop:nu-ratio-pdg}:
the $\approx 33.3$ prediction is no longer a postulate but follows from
explicit $\beta$-functions once D3 supplies $(M_{\mathrm{GUT}},M_{F,k})$.
The remaining open item is \emph{derivation} of those scales from
$V$ and observer-local FRG, not the two-loop counting.
\end{remark}

\begin{corollary}[D2 partial closure]\label{cor:flavon-d2}
Definition~\ref{def:flavon-potential}, Proposition~\ref{prop:flavon-beta},
Proposition~\ref{prop:flavon-beta-2loop}, Theorems~\ref{thm:nu-ratio-rg}
and~\ref{thm:nu-ratio-2loop}, and Proposition~\ref{prop:nu-ratio-pdg}
\textbf{partially close} development step~\textbf{D2}: the flavon potential,
one- and two-loop $\beta$-functions, factorized ratio
\eqref{eq:nu-ratio-factorized}, and PDG contact \eqref{eq:nu-ratio-pdg} are
explicit. Full closure requires D3 derivation of $M_{\mathrm{GUT}}$,
$M_{F,k}$, and $\mu_{\mathrm{osc}}$ without external scale input.
\end{corollary}